\section*{Capítulo \ref{ch:b2:blood-pressure} --- Regulación de la presión arterial y ejercicio}

\begin{solution}{exo:b2:blood-pressure:1}
$\bar P - P_{\text{ven}} = Q\,R = f\,V_s\,R$. La frecuencia cardíaca (el
\omterm{prop:b2:heart:conduction}{nódulo sinoauricular}, bajo el control del vago y de los nervios
simpáticos); el \omterm{prop:b2:heart:cycle}{volumen sistólico} (el \omterm{def:b2:heart:anatomy}{ventrículo}: su llenado, fijado por
las \omterm{def:b2:blood-circulation:vessels}{venas}, y su contractilidad, fijada por los nervios simpáticos); la
resistencia (el músculo liso de las \omterm{def:b2:blood-circulation:vessels}{arteriolas}, bajo el control de los
nervios simpáticos, las hormonas y los metabolitos locales).
\end{solution}

\begin{solution}{exo:b2:blood-pressure:2}
Sensores: receptores de estiramiento de los senos carotídeos y del cayado
aórtico. Centro: los centros cardiovasculares del bulbo raquídeo.
Efectores: el vago (frecuencia) y los nervios simpáticos (frecuencia,
contractilidad, tono arteriolar y venoso). Retroalimentación negativa,
con un retraso de alrededor de un segundo. Sin los sensores, la presión
media a lo largo de un día sigue siendo normal, pero la presión de un
minuto a otro oscila de forma desmesurada con cada cambio de postura y
cada emoción.
\end{solution}

\begin{solution}{exo:b2:blood-pressure:3}
El riñón (células de las \omterm{def:b2:blood-circulation:vessels}{arteriolas}) libera renina cuando caen su
perfusión o el aporte de sodio, o cuando se activan sus nervios
simpáticos; la renina corta el angiotensinógeno (hígado) en
angiotensina I; una enzima pulmonar la convierte en angiotensina II, que
contrae las \omterm{def:b2:blood-circulation:vessels}{arteriolas} y estimula la corteza suprarrenal para que segregue
aldosterona; la aldosterona hace que el riñón retenga sodio y agua, lo que
aumenta el volumen sanguíneo.
\end{solution}

\begin{solution}{exo:b2:blood-pressure:4}
La dilatación metabólica de las propias \omterm{def:b2:blood-circulation:vessels}{arteriolas} del músculo por sus
metabolitos; un \omterm{thm:b2:heart:starling}{gasto cardíaco} multiplicado por cinco, impulsado por los
nervios simpáticos y la bomba muscular; la constricción simpática de las
\omterm{def:b2:blood-circulation:vessels}{arteriolas} del intestino, los riñones y el músculo en reposo, que ceden
flujo.
\end{solution}

\begin{solution}{exo:b2:blood-pressure:5}
$83\times 1.2 = \qty{100}{mmHg}$; $333\times 0.35 = \qty{117}{mmHg}$.
La resistencia se divide por $1.2/0.35 = 3.4$: el radio se multiplica por $3.4^{1/4} =
1.36$.
\end{solution}

\begin{solution}{exo:b2:blood-pressure:6}
$30/(1 + G)$: 10, 6 y \qty{3.3}{mmHg}. Con la mitad de ganancia, la caída
residual se duplica: se marea al ponerse de pie, ve gris y puede
desmayarse.
\end{solution}

\begin{solution}{exo:b2:blood-pressure:7}
$\rho g h$: $1060\times 9.81\times 0.4 = \qty{4160}{Pa} =
\qty{31}{mmHg}$; $1060\times 9.81\times 1.2 = \qty{12500}{Pa} =
\qty{94}{mmHg}$. Encéfalo: $95 - 31 = \qty{64}{mmHg}$; tobillo: $95 + 94 =
\qty{189}{mmHg}$. Una cabeza de jirafa a \qty{2.5}{m} cuesta
\qty{195}{mmHg}: necesita una presión media de unos \qty{250}{mmHg} en el
\omterm{def:b2:heart:anatomy}{corazón}, y paredes vasculares gruesas y una piel tensa en las patas.
\end{solution}

\begin{solution}{exo:b2:blood-pressure:8}
$Q = 280/(190 - 130) = \qty{4.7}{L/min}$. Atleta: $5000/(200 - 30) =
\qty{29}{L/min}$; \omterm{prop:b2:heart:cycle}{volumen sistólico}, $\num{29400}/190 = \qty{155}{mL}$.
\end{solution}

\begin{solution}{exo:b2:blood-pressure:9}
$35/5 = \qty{7}{mmHg}$ menos: unos \qty{88}{mmHg}. Restablecen el volumen:
la reabsorción de líquido intersticial hacia los \omterm{def:b2:blood-circulation:vessels}{capilares} (de minutos a
una hora); la menor producción de orina bajo la \omterm{prop:b2:blood-pressure:raas}{hormona antidiurética} y la
aldosterona (de horas a días); la sed y la ingestión de agua (horas); los
\omterm{def:b2:blood-circulation:blood}{glóbulos rojos} de la médula ósea (semanas). La piel está pálida y fría
porque el reflejo ha cerrado sus \omterm{def:b2:blood-circulation:vessels}{arteriolas} para reservar la presión al
encéfalo y al \omterm{def:b2:heart:anatomy}{corazón}.
\end{solution}

\begin{solution}{exo:b2:blood-pressure:10}
Músculo con $R = 0.5$: $95/0.5 = \qty{190}{mL/s}$, \qty{11.4}{L/min};
con los \qty{4}{L/min} de los demás órganos, el \omterm{def:b2:heart:anatomy}{corazón} necesitaría
\qty{15.4}{L/min}. Con $Q$ fijo en \qty{5}{L/min}, la resistencia total
cae de 1,14 a $1/(1/1.14 - 1/5 + 1/0.5) =
\qty{0.37}{mmHg.s/mL}$ y la presión a \qty{31}{mmHg} --- el encéfalo
perdería el conocimiento. El organismo eleva el gasto hacia el primer caso
y contrae los demás lechos vasculares, de modo que la presión se mantiene
y el músculo recibe su flujo.
\end{solution}

\begin{solution}{exo:b2:blood-pressure:11}
(a) El riñón mal perfundido registró una presión baja y liberó renina, y
la angiotensina y la aldosterona elevaron la resistencia y el volumen.
(b) El \omterm{def:b2:blood-pressure:baroreflex}{barorreflejo} amortigua los cambios, pero se reajusta a cualquier
nivel que persista, y no puede modificar la retención de volumen del
riñón. (c) La fuente de renina, el órgano que exigía una presión más alta,
había desaparecido. (d) Habría bloqueado la conversión en angiotensina II:
menor resistencia, menos aldosterona, una presión más baja.
\end{solution}

\begin{solution}{exo:b2:blood-pressure:12}
El \omterm{def:b2:blood-pressure:baroreflex}{barorreflejo} actúa en un segundo, con una ganancia de alrededor de
cuatro y un valor de consigna que se desplaza hacia la presión dominante:
amortigua las fluctuaciones, pero no puede mantener un nivel. El riñón
actúa a lo largo de horas y días, con una ganancia en la práctica infinita
--- sigue excretando o reteniendo hasta alcanzar la presión a la que está
ajustado --- y su valor de consigna lo fija su propia fisiología: mantiene
el nivel. Una elevación crónica significa, por tanto, que el valor de
consigna del riñón se ha desplazado, y solo la reducen los fármacos que
actúan sobre el bucle renal o sobre la resistencia que exige.
\end{solution}

\begin{solution}{pb:b2:blood-pressure:1}
\textbf{1.} $1/R = 1/7.6 + 1/22.8 + 1/4.1 + 1/5.2 + 1/5.7 + 1/19 +
1/28.5 = 0.875$: $R = \qty{1.14}{mmHg.s/mL}$; $\bar P/Q = 95/83.3 =
1.14$.
\textbf{2.} $\bar P/R_i$: encéfalo 0,75; \omterm{def:b2:heart:anatomy}{corazón} 0,25; intestino e hígado
1,39; riñones 1,10; músculo 1,0; piel 0,30; otros \qty{0.20}{L/min};
fracciones de 15, 5, 28, 22, 20, 6 y \qty{4}{\%}.
\textbf{3.} $5000/70 = \qty{71}{mL}$; $5\times 50 = \qty{250}{mL}$ de
oxígeno por minuto.
\textbf{4.} Encéfalo: $0.75/1.4 = \qty{0.54}{L/min/kg}$; riñones: $1.1/0.3
= \qty{3.7}{L/min/kg}$, siete veces más que el encéfalo y cincuenta veces
la media del cuerpo: los riñones se irrigan para filtrar, no para
alimentarse.
\textbf{5.} Con una sola presión compartida por todos los órganos, cada
uno toma el flujo que fija su propia resistencia; regular los flujos de
forma centralizada privaría de \omterm{def:b2:blood-circulation:blood}{sangre} a cualquier órgano cuya demanda
cambiara. La presión es la moneda común.
\textbf{6.} \omterm{def:b2:heart:anatomy}{Corazón}: $250\times 0.2\times 0.7 = \qty{35}{mL/min}$;
músculo: $1000\times 0.2\times 0.25 = \qty{50}{mL/min}$. El \omterm{def:b2:heart:anatomy}{corazón}
extrae ya la mayor parte de lo que pasa, de modo que solo puede obtener
más oxígeno con más flujo, y por eso sus \omterm{def:b2:blood-circulation:vessels}{arteriolas} se dilatan en cuanto
trabaja más.
\textbf{7.} $Q$ baja un \qty{30}{\%} con $R$ fija: $\bar P$ baja un
\qty{30}{\%}, \qty{28.5}{mmHg}, hasta \qty{66.5}{mmHg}.
\textbf{8.} $28.5/5 = \qty{5.7}{mmHg}$: unos \qty{89}{mmHg}.
\textbf{9.} $Q = 85\times 0.7\times 71 = \qty{4.2}{L/min} =
\qty{70}{mL/s}$; $R = 89.3/70 = \qty{1.27}{mmHg.s/mL}$, un
\qty{11}{\%} más.
\textbf{10.} $1060\times 9.81\times 0.4 = \qty{4160}{Pa} =
\qty{31}{mmHg}$; encéfalo: $66.5 - 31 = \qty{35}{mmHg}$ antes del reflejo,
$89 - 31 = \qty{58}{mmHg}$ después.
\textbf{11.} $28.5/2 = \qty{14}{mmHg}$: \qty{81}{mmHg} en el \omterm{def:b2:heart:anatomy}{corazón},
\qty{50}{mmHg} en el encéfalo --- mareo, visión gris y casi un desmayo
cada vez que se levanta de una silla.
\textbf{12.} Tumbarse elimina la pérdida hidrostática de \qty{31}{mmHg}
y devuelve al \omterm{def:b2:heart:anatomy}{corazón} la \omterm{def:b2:blood-circulation:blood}{sangre} estancada; el \omterm{prop:b2:heart:cycle}{volumen sistólico} y la
presión en el encéfalo se recuperan en unos pocos latidos.
\textbf{13.} $45/5 = \qty{9}{mmHg}$: \qty{86}{mmHg}.
\textbf{14.} $1/R = 0.875 - 1/19 - 1/4.1 + 1/38 + 1/8.2 = 0.727$: $R =
\qty{1.38}{mmHg.s/mL}$. Encéfalo: $86/7.6 = \qty{11.3}{mL/s} =
\qty{0.68}{L/min}$; piel: $86/38 = \qty{2.3}{mL/s} = \qty{0.14}{L/min}$;
intestino: $86/8.2 = \qty{10.5}{mL/s} = \qty{0.63}{L/min}$.
\textbf{15.} Con menos \omterm{def:b2:blood-circulation:blood}{sangre} y \omterm{def:b2:blood-circulation:vessels}{arteriolas} más contraídas, la presión
\omterm{def:b2:blood-circulation:vessels}{capilar} cae por debajo de la \omterm{thm:b2:blood-circulation:oncotic}{presión oncótica} a lo largo de todo el
\omterm{def:b2:blood-circulation:vessels}{capilar}, de modo que el líquido se reabsorbe en lugar de filtrarse; el
\omterm{def:b2:blood-circulation:blood}{plasma} se diluye y el \omterm{def:b2:blood-circulation:blood}{hematocrito} baja.
\textbf{16.} La aldosterona, a través de la renina y la angiotensina,
desencadenada por la baja perfusión del riñón y la actividad simpática; la
\omterm{prop:b2:blood-pressure:raas}{hormona antidiurética}, desencadenada por la caída del volumen y el aumento
de la concentración plasmática.
\textbf{17.} $5\times 10^{12}$ células a $2.5\times 10^{6}$ por segundo:
$2\times 10^{6}$ s, unas tres semanas.
\textbf{18.} Por debajo de unos \qty{60}{mmHg} la curva del reflejo es
plana: los vasos están tan contraídos y el \omterm{def:b2:heart:anatomy}{corazón} late tan deprisa como
pueden, la ganancia es nula y cada nueva pérdida baja la presión en toda
su magnitud; los tejidos mal perfundidos liberan entonces ácido y se
dilatan, y la presión se desploma.
\textbf{19.} Músculo: $110/0.33 = \qty{333}{mL/s} = \qty{20}{L/min}$;
piel: $110/5 = \qty{22}{mL/s} = \qty{1.3}{L/min}$.
\textbf{20.} Resistencia del encéfalo, $8.8$ (flujo de 0,75 a 110); del
\omterm{def:b2:heart:anatomy}{corazón}, $6.6$
(flujo de 1,0): $1/R = 1/8.8 + 1/6.6 + 1/8.2 + 1/10.4 + 1/0.33 + 1/5 +
1/28.5 = 3.75$: $R = \qty{0.27}{mmHg.s/mL}$; $Q = 110/0.27 =
\qty{410}{mL/s} = \qty{25}{L/min}$.
\textbf{21.} $\num{24700}/180 = \qty{137}{mL}$.
\textbf{22.} $24.7\times(200 - 40) = \qty{3950}{mL/min}$, dieciséis veces
los \qty{250}{mL/min} de reposo.
\textbf{23.} Flujo, $\times 4.9$; extracción, $160/50 = \times 3.2$;
$4.9\times 3.2 = 16$.
\textbf{24.} La orden de la corteza motora y las señales de los receptores
de los músculos cambian el peso de la información de los \omterm{def:b2:blood-pressure:baroreflex}{barorreceptores}
en el bulbo raquídeo, de modo que el reflejo defiende \qty{110}{mmHg} en
lugar de 95 --- se eleva el valor de consigna, no se anula el bucle.
\textbf{25.} Caída residual al ponerse de pie, \qty{5.7}{mmHg};
\qty{86}{mmHg} tras la hemorragia; \qty{25}{L/min} y \qty{4}{L} de
oxígeno por minuto durante el ejercicio.
\end{solution}
